% SPDX-License-Identifier: CC-BY-4.0
\section{Evaluation}
\label{sec:evaluation}

\subsection{Setup}

\paragraph{Workloads.} Three synthetic factories drive the
experiments: \texttt{isotropic\_gaussian} for the canonical
experiments, \texttt{heterogeneous\_gaussian} for per-feature
heterogeneity, and \texttt{gaussian\_mixture} for multimodal benign
distributions. All are deterministic given a seed.

\paragraph{Attacks.} Three poisoning attacks instantiate the threat
model:
\begin{itemize}
\item \texttt{inject\_pre\_seed}, the black-box adversary,
  target-valued injections shuffled uniformly across the window.
\item \texttt{inject\_burst}, contiguous-block injections at a
  chosen or randomly-drawn position.
\item \texttt{inject\_targeted}, the white-box adversary that
  grid-searches the per-feature placement minimizing
  $\score(\target)$ after a fresh fit by an estimator with known
  \trim{}.
\end{itemize}

\paragraph{Estimators.} Four estimators are compared.
\textsc{Naive} uses per-feature mean and standard deviation, the
prior-work reference. \textsc{Consensus} splits the learning window
into four sub-windows and drops features whose sub-window means
disagree by more than three pooled standard deviations.
\textsc{Anomaly-Within} fits naive on the full window, drops the
top decile of within-window outliers, and refits on what survives.
\sediment{} runs at trim fractions $\{0.1, 0.2, 0.3, 0.5\}$.

\paragraph{Reporting.} Every cell uses ten independent seeds, drawn
deterministically from \texttt{numpy.random.SeedSequence(seed).spawn(4)}
to give independent learning-window, poisoning, benign-test, and
malicious-test streams. The target FPR is 0.05 throughout. Error
bars are standard error across seeds. Every figure regenerates from
\href{run:../scripts/reproduce_all.sh}{\texttt{scripts/reproduce\_all.sh}}.

\subsection{Clean-baseline FPR is at the configured target}
\label{sec:eval-clean-fpr}

By construction the per-estimator threshold is the
$1 - \texttt{target\_fpr}$ quantile of benign scores on each cell,
so the empirical FPR matches the configured target up to
discretization noise. \cref{fig:exp01} confirms this: every cell
lands at FPR\,=\,0.0500 with standard error below 0.0001 across ten
seeds. This is a sanity panel; differentiation across estimators
shows up in the next subsections.

\begin{figure}[t]
  \centering
  \includegraphics[width=\columnwidth]{fig01_clean_baseline_fpr.pdf}
  \caption{Clean-condition false-positive rate per estimator. Error
    bars are standard error across ten seeds.}
  \label{fig:exp01}
\end{figure}

\subsection{Naive collapse under poisoning}
\label{sec:eval-naive-collapse}

\cref{fig:exp02} shows the target score \textsc{Naive} assigns to
\target{} as the adversary's budget grows. Clean-condition score is
roughly 130. At $\budget = 0.05$ it falls to 70; at $\budget = 0.10$
to 48; at $\budget = 0.15$ to 36. The collapse is smooth and
approximately log-linear in \budget{}, not threshold-like. There is
no ``safe'' small budget. Even modest poisoning rates the adversary
could plausibly achieve cost substantial discriminability.

\begin{figure}[t]
  \centering
  \includegraphics[width=\columnwidth]{fig02_naive_under_poisoning.pdf}
  \caption{\textsc{Naive} target score under poisoning. The smooth
    decay shows there is no threshold-like ``safe'' adversary
    budget.}
  \label{fig:exp02}
\end{figure}

\subsection{\sediment{} design space}
\label{sec:eval-sediment-design}

\cref{fig:exp03} sweeps the joint $(\trim{}, \budget{})$ grid for
\sediment{}. Two structural features stand out. First, the
robustness boundary tracks $\budget = \trim/2$ predicted by
\cref{prop:robust}: for budgets at or below half the trim fraction,
the target score holds within a small constant factor of the
clean-condition value. Second, the cost of choosing too high a trim
is bounded: a $\trim = 0.50$ defender pays nothing in
clean-condition discriminability (\cref{sec:eval-fp-cost}) but
loses headroom against budgets above $0.25$. The recommended
operating point is the smallest $\trim$ at least double the worst
expected adversary budget.

\begin{figure}[t]
  \centering
  \includegraphics[width=\columnwidth]{fig03_robust_under_poisoning.pdf}
  \caption{\sediment{} target score across the budget $\times$
    trim-fraction grid. The robustness boundary tracks
    $\budget = \trim/2$.}
  \label{fig:exp03}
\end{figure}

\subsection{Budget sweep, naive vs.\ \sediment{}}
\label{sec:eval-budget-sweep}

\cref{fig:exp04-target} overlays \textsc{Naive} and \sediment{} at
$\trim = 0.30$ across the budget grid. At $\budget \leq 0.10$,
\sediment{}'s target score holds within 30\% of clean-condition
while \textsc{Naive} loses more than half. The gap closes past
$\budget = 0.20$ as \sediment{}'s budget guarantee runs out, but
the discriminability ratio stays favorable through $\budget = 0.30$.
The companion score-ratio and detection-rate panels in the
reproducibility artifact tell the same story with
threshold-invariant metrics.

\begin{figure}[t]
  \centering
  \includegraphics[width=\columnwidth]{fig04_target_score.pdf}
  \caption{Target score under poisoning, \textsc{Naive} versus
    \sediment{} at $\trim = 0.30$.}
  \label{fig:exp04-target}
\end{figure}

\subsection{Clean-condition FP cost of \trim{}}
\label{sec:eval-fp-cost}

The closed-form Gaussian-truncation correction $c(\trim)$ is meant
to restore the trimmed standard deviation to an unbiased estimate
of the underlying $\sigma$. \cref{fig:exp05} shows it works as
advertised. Across $\trim \in \{0.05, 0.10, 0.15, 0.20, 0.30, 0.40,
0.50\}$, both the threshold value and the empirical FPR are
statistically indistinguishable from the naive estimator's. Under
Gaussian benign distributions, picking a higher trim costs nothing
in clean-condition false-positive rate. The cost shows up only
when the benign distribution departs from the Gaussian null, which
we look at next.

\begin{figure}[t]
  \centering
  \includegraphics[width=\columnwidth]{fig05_fp_cost.pdf}
  \caption{Clean-condition FPR as a function of trim fraction. Under
    Gaussian benign, the Gaussian-truncation correction holds the
    cost to within statistical noise.}
  \label{fig:exp05}
\end{figure}

\subsection{Multimodal benign distributions}
\label{sec:eval-multimodal}

\cref{fig:exp06} runs the budget sweep against a two-component
Gaussian mixture benign with the target placed past the upper mode.
The closed-form correction $c(\trim)$ is calibrated to a unimodal
Gaussian null and overestimates how much variance the trimmed tails
were carrying, so the corrected $\sigma$ underestimates the true
population spread. The visible effect is inflated absolute target
scores at low budgets compared with the unimodal case. The
threshold-invariant discriminability ratios stay reasonable. We
return to this structural limit in \cref{sec:discussion}.

\begin{figure}[t]
  \centering
  \includegraphics[width=\columnwidth]{fig06_target_score.pdf}
  \caption{Target score under multimodal benign distributions.}
  \label{fig:exp06}
\end{figure}

\subsection{Ablation: alternative robustness primitives, uniform poisoning}
\label{sec:eval-ablation}

\cref{fig:exp07} runs the four-way ablation under \texttt{pre\_seed}
poisoning. \sediment{} dominates at every budget tested.
\textsc{Anomaly-Within} provides meaningful resistance up to roughly
its own \texttt{drop\_fraction} (0.10) and degrades smoothly past
it. \textsc{Consensus} is mathematically identical to
\textsc{Naive} under uniform poisoning. \texttt{pre\_seed} shuffles
injections uniformly across sub-windows, so every sub-window's mean
shifts by the same amount, the disagreement test never fires, and
the score reduces to the unfiltered naive computation.

\begin{figure}[t]
  \centering
  \includegraphics[width=\columnwidth]{fig07_target_score.pdf}
  \caption{Four-way ablation under uniform poisoning.}
  \label{fig:exp07}
\end{figure}

\subsection{Ablation: alternative robustness primitives, burst poisoning}
\label{sec:eval-burst}

\cref{fig:exp08} repeats the four-way ablation under
\texttt{inject\_burst}, the attack pattern that should give a
sub-window consensus filter the strongest chance. The result runs
counter to that intuition: \textsc{Consensus} performs \emph{worse}
than \textsc{Naive} at high burst budgets. The features carrying
the largest target signal are also the features whose sub-window
means disagree most under a burst, so the consensus filter drops
exactly the features that would have driven detection. The
ablation surfaces a structural problem with consensus filtering as
a robustness primitive against this threat model: it cannot tell a
feature shifted by an attacker burst apart from a feature shifted
by an underlying workload change, and penalizing the former drops
detection signal along with it.

\begin{figure}[t]
  \centering
  \includegraphics[width=\columnwidth]{fig08_target_score.pdf}
  \caption{Four-way ablation under burst poisoning.}
  \label{fig:exp08}
\end{figure}

\subsection{White-box adversary}
\label{sec:eval-white-box}

\cref{fig:exp09} repeats the four-way ablation under
\texttt{inject\_targeted}, the grid-search white-box adversary that
knows each defender's trim parameter. The headline finding: the
optimum placement against \sediment{} converges to the at-target
injection. We verified this numerically by inspecting the
grid-search trajectory. At the budgets where the optimum is
interior to $[\text{median}, \target]$, the chosen value sits at
the trim-boundary benign quantile, and the resulting fitted moments
are mathematically equivalent to those produced by at-target
injection. Swapping one boundary-quantile benign sample for one
boundary-valued injection leaves the trimmed mean and variance
unchanged. White-box knowledge of \trim{} therefore does not yield
a stronger attack against a symmetric-trim defender than the
budget-bounded black-box adversary the construction was designed
against. This is the strongest threat-model evidence the paper
produces. We discuss residual risks from more sophisticated
adaptive adversaries in \cref{sec:discussion}.

\begin{figure}[t]
  \centering
  \includegraphics[width=\columnwidth]{fig09_target_score.pdf}
  \caption{Four-way ablation under white-box targeted poisoning.
    The grid-search adversary's optimal placement against
    \sediment{} reduces to the at-target injection.}
  \label{fig:exp09}
\end{figure}

\subsection{Real-data validation}
\label{sec:eval-real-data}

The synthetic experiments above are the paper's controlled evidence.
This subsection reports a smaller-scale end-to-end validation against
real Linux behavior captured live by the runtime, included to
demonstrate that the algorithm transfers to a real captured feature
distribution.

\paragraph{Setup.} Ubuntu 24.04 (kernel 6.6) under WSL2. The
\chronosynd{} BPF probe attaches to \texttt{raw\_syscalls/sys\_enter}
and emits one event per syscall, classified by syscall number against
a 26-entry vocab plus an other-bucket. The collector ran for 3
minutes capturing real \texttt{bash} activity (a script driving
\texttt{ls}, \texttt{cat}, \texttt{find}, \texttt{grep}, \texttt{wc}
in a loop), producing 730k bash-keyed events forming 45{,}640
closed feature windows. \sediment{} fitted at $\trim = 0.30$ on the
captured trace. The persisted baseline showed nonzero mass at six
vocab buckets (read, close, clone, execve, openat, other) with
matching nonzero standard deviations, confirming a non-degenerate
fit on real workload data.

\paragraph{Attack run.} A separate 60-second capture ran the
\texttt{spawn\_storm.sh} (200 helper processes) and
\texttt{etc\_recon.sh} (500 file reads under \texttt{/etc}) payloads
under bash while the daemon scored bash windows live against the
fitted baseline. Both payloads shape bash's syscall mix sharply
toward exec/clone/openat events.

\paragraph{Results.} The fitted threshold for bash was 100.0. Across
3{,}846 bash windows during the attack capture:

\begin{itemize}
\item \textbf{Alerts:} 1{,}683 windows scored above threshold (43.8\%).
\item \textbf{Clean (below threshold):} 2{,}163 windows.
\item \textbf{Clean window score distribution:} median 4.83, p90 30.69,
  max 83.01. Every clean window scored well below the 100 threshold,
  so the empirical false-positive rate during the attack window was
  zero.
\item \textbf{Alert window score distribution:} median 609.51, p90
  $7.8 \times 10^{9}$, max $3.2 \times 10^{11}$.
\item \textbf{Median separation:} the median alert score is 126$\times$
  the median clean score.
\end{itemize}

The very large alert scores are real but reflect a known artifact of
the fitted baseline: when a feature has zero variance under clean
bash (e.g., a syscall that never occurs in the clean window), the
trimmed std collapses to zero, the per-feature score denominator
floors at $\epsilon$, and any nonzero attack value blows the squared
standardized residual up by $1/\epsilon^2$. The construction
correctly flags the deviation, the magnitude just compresses
poorly into a single scalar at the extreme. A logarithmic score
or a per-feature cap at three or four standard deviations would
give cleaner numerics; we leave that as a deployment-time tuning
choice rather than an algorithmic change.

This is a single workload on a single host, not a controlled
multi-host evaluation. We report it as evidence the artifact runs
end-to-end on real captured behavior with non-trivial detection
signal, not as a substitute for the controlled synthetic
experiments above.
